\section{Discusión}

Aunque los resultados numéricos lucen prometedores sobre el papel, su interpretación real exige mirar más allá de las métricas. Resulta revelador observar cómo la brecha entre simulación y operación diaria suele decidir el destino de proyectos satelitales en entornos complejos.

\subsection{Limitaciones técnicas y operativas}

Uno de los desafíos persistentes que surge al analizar arquitecturas redundantes radica en las limitaciones inherentes que persisten incluso después de implementar las mejores prácticas. La configuración propuesta mejora sustancialmente la disponibilidad, pero no elimina riesgos correlacionados como tormentas solares intensas o ciberataques sofisticados que podrían afectar simultáneamente rutas primaria y de respaldo.

En muchos casos estudiados, la dependencia de componentes importados y la escasez de talento técnico altamente especializado local constituyen cuellos de botella operativos reales. Además, la banda Ka, aunque atractiva por su capacidad, impone requisitos de pointing extremadamente precisos que pueden verse comprometidos durante eventos meteorológicos severos. 

Trabajos sobre sostenibilidad en comunicaciones satelitales destacan que estas limitaciones no solo son técnicas, sino que impactan directamente la huella de carbono y los costos de ciclo de vida del sistema \cite{SustainableSat2022}. En el contexto venezolano, factores como la disponibilidad de energía estable en estaciones remotas y la logística de mantenimiento en zonas de difícil acceso añaden capas adicionales de complejidad que no siempre se capturan plenamente en simulaciones.

\subsection{Impacto económico y social}

Particularmente llamativo es el hecho de que la inversión en redundancia satelital trasciende con creces el ámbito técnico. Una disponibilidad superior al 99.95\% puede traducirse en mejoras tangibles para servicios educativos, telemedicina y respuesta a emergencias en regiones aisladas del país. 

Sin embargo, el costo incremental de la arquitectura redundante exige una evaluación cuidadosa. Aunque los beneficios operativos justifican la inversión a mediano plazo, las restricciones presupuestarias actuales obligan a considerar implementaciones por fases. El impacto social potencial —reducción de la brecha digital y fortalecimiento de la soberanía en comunicaciones— parece sugerir que los retornos no económicos bien pueden compensar parte de la carga financiera.

No obstante, cabe matizar que estos beneficios solo se materializarán si se acompaña el despliegue con programas de capacitación local y desarrollo de industria auxiliar. De lo contrario, el sistema podría terminar dependiendo excesivamente de soporte externo, limitando su sostenibilidad a largo plazo.

\subsection{Recomendaciones}

Lejos de ser un asunto resuelto, el camino hacia la implementación efectiva del VENESAT-2 redundante requiere decisiones estratégicas claras. En primer lugar, se recomienda priorizar la diversidad geográfica de estaciones terrenas y la adquisición progresiva de capacidad de respaldo en otros operadores GEO durante la fase inicial.

Se sugiere además establecer un programa de monitoreo continuo de propagación para refinar los modelos locales y validar las simulaciones aquí presentadas. La integración de técnicas emergentes de inteligencia artificial para predicción de atenuación y enrutamiento dinámico podría incorporarse en una segunda etapa.

Mirando hacia el futuro, resulta estratégico alinear este proyecto con las tendencias de evolución hacia sistemas 6G, donde la convergencia satelital-terrestre y la multi-conectividad jugarán roles centrales \cite{SatComm6G2025}. Recomendamos específicamente destinar recursos a la formación de ingenieros locales en tecnologías satelitales avanzadas y explorar modelos de asociación público-privada que mitiguen riesgos financieros.

En síntesis, los hallazgos de este estudio indican que la implementación de un sistema redundante para VENESAT-2 es técnicamente viable y operativamente deseable. Su éxito dependerá, en gran medida, de la capacidad para adaptar las soluciones técnicas a las realidades locales sin perder de vista los objetivos estratégicos de soberanía y desarrollo inclusivo.